% !TEX TS-program = pdflatex
% !TeX program = pdflatex
% !TEX encoding = UTF-8
% !TEX spellcheck = en_US

\documentclass[xcolor=table]{beamer}

\usepackage{../extra/beamer/karimnlp}

\input{options}

\subtitle[05- POS tagging]{Chapter 05\\Part-of-speech tagging} 

\changegraphpath{../img/pos/}

\begin{document}

\begin{frame}
	\frametitle{\inserttitle}
	\framesubtitle{\insertshortsubtitle: Introduction}
	
	\begin{exampleblock}{Ambiguous sentences}
		\begin{center}
			\Large\bfseries
			We can can the can.\\[0.3em]
			Will Will will the will to Will?
		\end{center}
	\end{exampleblock}
	
	\begin{itemize}
%		\item These sentences illustrate lexical and syntactic ambiguity.
		\item Identify the part of speech of each word.
		\item \optword{can}:
		\begin{itemize}
			\item Verb (modal): express ability
			\item Verb (lexical): preserve/put into a container
			\item Noun: container
		\end{itemize}
		
		\item \optword{will}:
		\begin{itemize}
			\item Modal auxiliary verb: express future
			\item Lexical verb: intend or wish
			\item Noun: testament
			\item Proper noun: a person's name
		\end{itemize}
	\end{itemize}
	
\end{frame}


\begin{frame}
	\frametitle{\inserttitle}
	\framesubtitle{\insertshortsubtitle: Some humor}

	\begin{center}
		\vgraphpage{humor/humor-pos.jpg}
	\end{center}

\end{frame}

\begin{frame}
	\frametitle{\inserttitle}
	\framesubtitle{\insertshortsubtitle: Plan}

	\begin{multicols}{2}
	%	\small
	\tableofcontents
	\end{multicols}

\end{frame}

%===================================================================================
\section{Sequence labeling}
%===================================================================================

\begin{frame}
	\frametitle{\insertshortsubtitle}
	\framesubtitle{\insertsection}
	
	\begin{itemize}
		\item In most cases, we consider a sequence of words.
		\item Each word is assigned a label (or tag).
		\item This label does not depend solely on the word's features.
		\item It also depends on the labels of neighboring words.
		\item Some tasks assign labels to spans (chunks) of words rather than individual tokens.
		\item In such cases, the \keyword{IOB} notation (introduced in \textbf{Chapter~02}) is commonly used.
	\end{itemize}
	
\end{frame}

\subsection{Description}

\begin{frame}
	\frametitle{\insertshortsubtitle: \insertsection}
	\framesubtitle{\insertsubsection}

	\begin{itemize}
		\item \textbf{Goal}: Assigning labels to the elements of a sequence.
		\item \textbf{Example}: \expword{Assigning POS tags to the words in a sentence}.
	\end{itemize}
	
	\begin{block}{Problem formulation}
		\begin{itemize}
			\item $w = (w_1, \ldots, w_n)$: input word sequence
			\item $t = (t_1, \ldots, t_n)$: output label sequence
		\end{itemize}
		\begin{center}
			$\arg\max\limits_t P(t \mid w)$
		\end{center}
		
		\begin{itemize}
			\item \optword{Generative models}: 
			$\arg\max\limits_t P(t \mid w) = \arg\max\limits_t P(t)\,P(w \mid t)$ % \quad (since $P(w)$ is constant)
			\item \optword{Discriminative models}: directly estimate $\arg\max\limits_t P(t \mid w)$
		\end{itemize}
	\end{block}

\end{frame}

\subsection{Applications}

\begin{frame}
	\frametitle{\insertshortsubtitle: \insertsection}
	\framesubtitle{\insertsubsection}

	\begin{itemize}
		\item Named Entity Recognition (NER)
		\begin{itemize}
			\item Identify entities such as persons, organizations, and locations in a sentence
		\end{itemize}
		
		\item POS Tagging
		\begin{itemize}
			\item Assign part-of-speech (POS) tags (e.g., noun, verb) to each word
		\end{itemize}
		
		\item Chunking
		\begin{itemize}
			\item Identify syntactic phrases (e.g., noun phrases) in a sentence
		\end{itemize}
	\end{itemize}

\end{frame}

\begin{frame}
	\frametitle{\insertshortsubtitle: \insertsection}
	\framesubtitle{\insertsubsection: Example}

	\begin{figure}
		\centering
		\hgraphpage{exp-ner2_.pdf}
		\caption{Example of NER [\url{https://corenlp.run/}]}
	\end{figure}
	
	\begin{figure}
		\centering
		\hgraphpage{exp-pos2_.pdf}
		\caption{Example of POS tagging [\url{https://corenlp.run/}]}
	\end{figure}

\end{frame}


%===================================================================================
\section{Task description}
%===================================================================================

\begin{frame}
\frametitle{Part-of-speech tagging}
\framesubtitle{Task description}

	\begin{itemize}
		\item \textbf{Objective}: assign a part-of-speech (POS) tag to each word in a sentence
		
		\item POS tag sets:
		\begin{itemize}
			\item General: \expword{NOUN, VERB, ADJ, ADV, ADP, DET}
			\item Specific: \expword{NN (noun), NNP (proper noun), VB (base verb), VBD (past tense), ...}
		\end{itemize}
		
		\item Corpus annotation:
		\begin{itemize}
			\item A test corpus is manually annotated.
			\item A separate corpus is used for training in supervised learning settings.
		\end{itemize}
	\end{itemize}

\end{frame}

\subsection{Universal classes}

\begin{frame}
	\frametitle{\insertshortsubtitle: \insertsection}
	\framesubtitle{\insertsubsection}

	\url{https://universaldependencies.org/u/pos/}
	
	\begin{tblr}{
			colspec = {p{.3\textwidth}p{.32\textwidth}p{.28\textwidth}},
			row{odd} = {lightblue},
			row{even} = {lightyellow},
			row{1} = {darkblue},
			rowsep = 1pt,
		} 
		\textcolor{white}{Open class} & \textcolor{white}{Close class} & \textcolor{white}{Other} \\
		
		\keyword{ADJ}:  adjective & \keyword{ADP}: adposition & \keyword{PUNCT}: punctuation \\
		\keyword{ADV}:  adverb & \keyword{AUX}: auxiliary & \keyword{SYM}: symbol \\
		\keyword{INTJ}: interjection & \keyword{CCONJ}: coordination conjunction & \keyword{X}: Other \\
		\keyword{NOUN}: noun & \keyword{DET}: determiner &  \\
		\keyword{PROPN}: proper noun & \keyword{NUM}: numerical &  \\
		\keyword{VERB}: verb & \keyword{PART}: particle &  \\
		 & \keyword{PRON}: pronoun &  \\
		 & \keyword{SCONJ}: subordination conjunction &  \\
		
	\end{tblr}

\end{frame}

\subsection{Treebanks}

\begin{frame}[fragile]
	\frametitle{\insertshortsubtitle: \insertsection}
	\framesubtitle{\insertsubsection}
	
	\vspace{-8pt}
	\begin{figure}
		\begin{tcolorbox}[colback=white, colframe=blue, boxrule=1pt, text width=.9\textwidth]
			\small
	\begin{alltt}
	Battle-tested\keyword{/JJ} Japanese\keyword{/JJ} industrial\keyword{/JJ} managers\keyword{/NNS}
	here\keyword{/RB} always\keyword{/RB} buck\keyword{/VBP} up\keyword{/RP} nervous\keyword{/JJ} newcomers\keyword{/NNS}
	with\keyword{/IN} the\keyword{/DT} tale\keyword{/NN} of\keyword{/IN} the\keyword{/DT} first\keyword{/JJ} of\keyword{/IN}
	their\keyword{/PP\$} countrymen\keyword{/NNS} to\keyword{/TO} visit\keyword{/VB} Mexico\keyword{/NNP} ,\keyword{/,}
	a\keyword{/DT} boatload\keyword{/NN} of\keyword{/IN} samurai\keyword{/FW} warriors\keyword{/NNS} blown\keyword{/VBN}
	ashore\keyword{/RB} 375\keyword{/CD} years\keyword{/NNS} ago\keyword{/RB} .\keyword{/.}
	\end{alltt}\vspace{-6pt}
	\end{tcolorbox}
	\caption{Example from Penn Treebank \cite{2003-taylor}}
	\end{figure}
	
	\begin{itemize}
		\item Universal Treebank (\url{https://universaldependencies.org/}) \cite{2012-petrov-al}
		\begin{itemize}
			\item It is an open community project.
			\item It aims to provide annotated corpora for several languages.
			\item It uses the same PoS for all languages.
		\end{itemize}
	\end{itemize}

\end{frame}

\subsection{Difficulties and tools}

\begin{frame}
	\frametitle{\insertshortsubtitle: \insertsection}
	\framesubtitle{\insertsubsection: Difficulties}

	\begin{itemize}
		\item \optword{Ambiguity}
		\begin{itemize}
			\item Some words can have multiple POS tags.
			\item E.g., \expword{We can can the can.}
		\end{itemize}
		
		\item \optword{Unknown words}
		\begin{itemize}
			\item Out-of-vocabulary (OOV) words (i.e., not seen in the training corpus).
			\item Languages are constantly evolving, leading to the creation of new words and abbreviations.
		\end{itemize}
		
		\item \optword{Tag inconsistency}
		\begin{itemize}
			\item Differences in tag sets may lead to inconsistencies in annotation
			\item E.g., \expword{In the ``Brown'' and ``WSJ'' corpora, the word ``to'' is tagged as ``TO''. In the ``Switchboard'' corpus, it is tagged as ``TO'' when marking the infinitive, and as ``IN'' when used as a preposition.}
		\end{itemize}
	\end{itemize}

\end{frame}

\begin{frame}
	\frametitle{\insertshortsubtitle: \insertsection}
	\framesubtitle{\insertsubsection: Tools}

%	\begin{itemize}
%		\item \url{https://nlp.stanford.edu/software/tagger.shtml}
%		\item \url{https://github.com/sloria/textblob}
%		\item \url{https://www.nltk.org/api/nltk.tag.html}
%		\item \url{https://spacy.io/}
%		\item \url{https://github.com/clips/pattern}
%	\end{itemize}

	\begin{itemize}
		\item \textbf{Stanford POS Tagger}
		\begin{itemize}
			\item \url{https://nlp.stanford.edu/software/tagger.shtml}
			\item A classical probabilistic POS tagging system
		\end{itemize}
		
		\item \textbf{NLTK}
		\begin{itemize}
			\item \url{https://www.nltk.org/api/nltk.tag.html}
			\item A widely used educational NLP toolkit in Python
		\end{itemize}
		
		\item \textbf{spaCy}
		\begin{itemize}
			\item \url{https://spacy.io/}
			\item A modern, efficient library for industrial NLP applications
		\end{itemize}
		
		\item \textbf{TextBlob}
		\begin{itemize}
			\item \url{https://github.com/sloria/textblob}
			\item A simple NLP library built on top of NLTK
		\end{itemize}
		
		\item \textbf{Pattern}
		\begin{itemize}
			\item \url{https://github.com/clips/pattern}
			\item A lightweight NLP library for Python
		\end{itemize}
		
	\end{itemize}

\end{frame}

\begin{frame}
	\frametitle{\insertshortsubtitle: \insertsection}
	\framesubtitle{\insertsubsection: Some humor}

	\begin{center}
		\vgraphpage{humor/humor-identify.jpg}
	\end{center}

\end{frame}

%===================================================================================
\section{Approaches}
%===================================================================================

\begin{frame}
\frametitle{Part-of-speech tagging}
\framesubtitle{Approaches}

	\begin{itemize}
		\item Rule-based approaches
		
		\item Machine learning approaches
		\begin{itemize}
			\item Transformation-Based Learning (TBL)
			\item Statistical models: \expword{Hidden Markov Models (HMMs), Maximum Entropy Markov Models (MEMMs), Conditional Random Fields (CRFs)}
			\item Neural models: \expword{Recurrent Neural Networks (RNNs), Transformers}
		\end{itemize}
	\end{itemize}

\end{frame}

\subsection{Hidden Markov Model (HMM)}

\begin{frame}[fragile]
	\frametitle{\insertshortsubtitle: \insertsection}
	\framesubtitle{\insertsubsection}
	
	\begin{itemize}
		\item $Q = \{t_1, t_2, \ldots, t_N\}$: set of tags (states)
		\item $A$: transition probability matrix, where $\sum_j a_{ij} = 1,\ \forall i$
		\item $W = (w_1, w_2, \ldots, w_T)$: word sequence
		\item $B$: emission probabilities $b_i(w_t)$; probability of generating word $w_t$ in state $t_i \in Q$
		\item $\pi = (\pi_1, \pi_2, \ldots, \pi_N)$: initial state distribution, where $\sum_i \pi_i = 1$
		
		\item $\hat{t} = \arg\max\limits_t P(t \mid w) = \arg\max\limits_t P(t)\,P(w \mid t)$
		
		\item $P(t) = P(t_1, \ldots, t_T) \approx \pi_{t_1} \prod\limits_{i=2}^{T} P(t_i \mid t_{i-1})$ \quad (bigram assumption)
		
		\item $P(w \mid t) = \prod\limits_{i=1}^{T} P(w_i \mid t_i)$
	\end{itemize}

\end{frame}

\begin{frame}
	\frametitle{\insertshortsubtitle: \insertsection}
	\framesubtitle{\insertsubsection: Example}

	\begin{exampleblock}{Example of a training corpus}
		\begin{itemize}
			\item a/DT computer/NN can/VB help/VB you/PN
			\item he/PN wants/VB to/TO help/VB you/PN
			\item he/PN wants/VB a/DT computer/NN
			\item he/PN can/VB swim/VB
		\end{itemize}
	\end{exampleblock}
	
	\begin{itemize}
		\item $P(VB | PN) = \frac{C(PN,\ VB)}{C(PN)} = \frac{3}{5}$, 
			  $P(VB | VB) = \frac{C(VB,\ VB)}{C(VB)} = \frac{2}{7}$,
		\item $P(he | PN) = \frac{C(PN,\ he)}{C(PN)} = \frac{3}{5}$,
		      $P(can | VB) = \frac{C(VB,\ can)}{C(VB)} = \frac{2}{7}$,
			  $P(help | VB) = \frac{C(VB,\ help)}{C(VB)} = \frac{2}{7}$,
		\item $\pi_{PN} = \frac{C(PN,\ <s>)}{C(<s>)} = \frac{3}{4} $
		\item $P(PN\ VB\ VB | he\ can\ help) \approx (\pi_{PN} P(VB | PN) P(VB | VB)) (P(he | PN) P(can | VB) P(help | VB)) $
	\end{itemize}

\end{frame}

%\begin{frame}
%	\frametitle{\insertshortsubtitle: \insertsection}
%	\framesubtitle{\insertsubsection: Exercise}
%	
%	\begin{itemize}
%		\item Use the example training corpus from the previous slide.
%		\item Estimate HMM parameters: transition matrix $A$, emission matrix $B$, and initial distribution $\pi$.
%		\item Calculate the approximate probability $P(PN\ VB\ VB \mid he\ can\ help)$.
%		\item Using greedy decoding, determine the tag sequence for the sentence \expword{he can help}.
%		\item Repeat the procedure using Viterbi decoding.
%		\item Repeat using Beam Search decoding with $K=3$.
%	\end{itemize}
%	
%\end{frame}

\subsection{Maximum Entropy Markov Model}

\begin{frame}
	\frametitle{\insertshortsubtitle: \insertsection}
	\framesubtitle{\insertsubsection}

	\begin{itemize}
		\item We denote a subsequence: $x_n^m = (x_n, x_{n+1}, \ldots, x_m)$.
		\item Given:
		\begin{itemize}
			\item a word sequence: $w = (w_1, w_2, \ldots, w_n)$
			\item a set of feature functions $f_j$ defined on a window of words $w_{i-l}^{i+l}$ and the previous $k$ tags $t_{i-k}^{i-1}$,
		\end{itemize}
		\item tag sequence prediction:
		\[
		\hat{t} = \arg\max_{t} P(t \mid w) \approx \arg\max_{t} \prod_{i=1}^{n} P(t_i \mid w_{i-l}^{i+l}, t_{i-k}^{i-1})
		\]
		\item Conditional probability (softmax over features):
		\[
		P(t_i \mid w_{i-l}^{i+l}, t_{i-k}^{i-1}) =
		\frac{\exp\left(\sum_j \theta_j f_j(t_i, w_{i-l}^{i+l}, t_{i-k}^{i-1})\right)}
		{\sum_{t'_i} \exp\left(\sum_j \theta_j f_j(t'_i, w_{i-l}^{i+l}, t_{i-k}^{i-1})\right)}
		\]
	\end{itemize}

\end{frame}

\begin{frame}
	\frametitle{\insertshortsubtitle: \insertsection}
	\framesubtitle{\insertsubsection: Features}

	\begin{itemize}
		\item Current word $w_i$ (words themselves as features)
		\item Previous tags $t_{i-1}, t_{i-2}, \dots$
		\item $w_i$ contains a prefix from a predefined list ($\text{length} \le 4$)
		\item $w_i$ contains a suffix from a predefined list ($\text{length} \le 4$)
		\item $w_i$ contains a digit
		\item $w_i$ contains an uppercase letter
		\item $w_i$ contains a hyphen
		\item $w_i$ is entirely uppercase
		\item $w_i$ template (e.g., \expword{X.X.X})
		\item $w_i$ is uppercase with a hyphen and a number (e.g., \expword{CFC-12})
		\item $w_i$ is uppercase followed by an abbreviation (e.g., \expword{Co., Inc.}) within a window of up to 3 words
	\end{itemize}

\end{frame}

%\begin{frame}
%	\frametitle{\insertshortsubtitle: \insertsection}
%	\framesubtitle{\insertsubsection: Exercise}
%	
%	\begin{itemize}
%		\item Use the training corpus from the previous HMM example.
%		\item We will use the following features:
%		\begin{itemize}
%			\item Past, current, and future words, encoded using ordinal encoding (alphabetical order).
%			\item Past and current PoS tags, encoded using ordinal encoding (alphabetical order).
%			\item Indicator feature: current word ends with the suffix "s" (1 if yes, 0 otherwise).
%		\end{itemize}
%		\item Manually train a MEMM model with the following specifications:
%		\begin{itemize}
%			\item Parameters are initialized to 1; no bias term is used.
%			\item Use \textbf{average normalization} instead of Softmax as the activation function.
%			\item Replace cross-entropy loss with $J = -\sum_{i=1}^{M} t_i \hat{t}_i$.
%			\item Learning rate $\alpha = 1$.
%			\item Number of iterations: 3.
%		\end{itemize}
%	\end{itemize}
%	
%\end{frame}

\subsection{Neural model}

\begin{frame}
	\frametitle{\insertshortsubtitle: \insertsection}
	\framesubtitle{\insertsubsection: Simple RNN}

	\hgraphpage{rnn-simple.pdf}

\end{frame}

\begin{frame}
	\frametitle{\insertshortsubtitle: \insertsection}
	\framesubtitle{\insertsubsection: Stacked RNN}

	\hgraphpage{rnn-stack.pdf}

\end{frame}

\begin{frame}
	\frametitle{\insertshortsubtitle: \insertsection}
	\framesubtitle{\insertsubsection: Bidirectional RNN}

	\hgraphpage{rnn-bi.pdf}

\end{frame}

\begin{frame}
	\frametitle{\insertshortsubtitle: \insertsection}
	\framesubtitle{\insertsubsection: Character-based embedding RNN}

	\hgraphpage{rnn-char.pdf}

\end{frame}

\begin{frame}
	\frametitle{\insertshortsubtitle: \insertsection}
	\framesubtitle{\insertsubsection: Transformers}
	
	\hgraphpage{transformers.pdf}
	
\end{frame}

\begin{frame}
	\frametitle{\insertshortsubtitle: \insertsection}
	\framesubtitle{\insertsubsection: Some humor}

	\begin{center}
		\vgraphpage{humor/humor-learn.jpg}
	\end{center}

\end{frame}

\insertbibliography{NLP05}{*}

\end{document}

